\section{Nonlinear arithmetic: exact certificates, not floating-point proofs}
\label{sec:nonlinear}

\subsection{The acceptance language}

The nonlinear worker targets polynomial statements over $\R$ with exact rational
coefficients. Variables may stand for opaque real subexpressions; any further
relationship between them must be supplied as a proved constraint. The same
representation must not be reused over an arbitrary ring with an assumed real
ordering.

Normalise the target to $p(\bm x)\ge0$, the hypotheses to $g_i(\bm x)\ge0$, and
the equalities to $h_j(\bm x)=0$. The certificate language is
\begin{equation}
 p=\sum_{k=1}^{s}c_k\,q_k^2\prod_{i=1}^{m}g_i^{a_{ki}}
        +\sum_{j=1}^{r}u_j\,h_j,
 \qquad c_k\in\Q_{\ge0},\quad a_{ki}\in\N,
 \label{eq:cone}
\end{equation}
with $q_k,u_j$ rational polynomials. The checker is given the original
$p,g_i,h_j$ independently of the certificate, reconstructs the right-hand side
exactly, and requires polynomial identity.

\begin{proposition}[Cone-certificate soundness]
\label{prop:cone}
If \eqref{eq:cone} holds, every point satisfying $g_i\ge0$ and $h_j=0$ satisfies
$p\ge0$.
\end{proposition}
\begin{proof}
Each square is nonnegative, each hypothesis product is nonnegative, and each
coefficient $c_k$ is nonnegative. Every equality-multiplier term vanishes.
Summing proves the claim.
\end{proof}

This elementary proof is exactly why search can be far more complicated than
checking --- the observation Harrison made foundational for SOS
reconstruction~\cite{harrison-sos}. The language is a bounded preordering-style
cone with equality-ideal terms. Restricting products to individual $g_i$ gives a
smaller family; fixing the possible $q_k$ gives a finite cone. None of these
should be called an unrestricted positivity decision procedure.

Two details of \eqref{eq:cone} matter in implementation. The equality
multipliers $u_j$ carry \emph{no sign condition}: an LP formulation with
nonnegative variables therefore needs $\pm$ columns for each ideal generator.
The reason is not pedantry --- under $y=x$,
\[
 y^2-2x+1=(x-1)^2+(x+y)(y-x),
\]
and the second term need not be nonnegative. It vanishes because of the
equality. Treating every multiplier as nonnegative would discard that
flexibility. The second detail is naming: in the serialised formats the field
called \code{square} stores $q_k$, not $q_k^2$.

\paragraph{Strict inequalities.}
Strictness needs explicit information. One sufficient certificate is an identity
of the same form for $p-\varepsilon$ with rational $\varepsilon>0$. It is not
necessary in general: a polynomial can be strictly positive on an unbounded
domain while having infimum zero, so $x>0$ implies $x>0$ without a uniform
rational lower bound. Another route is a term whose coefficient and factors are
provably strictly positive at every admissible point. The proof layer must
preserve the strict inequality rather than silently weakening it to
nonnegativity.

\paragraph{Refutation.}
A certificate for $-1$ in the nonnegative cone plus the equality ideal gives a
contradiction. If $q>0$ is a strict hypothesis, a checked identity expressing
$-q^m$ as a cone element plus vanishing ideal terms, with $m\ge1$, also gives
one.

\subsection{Search is untrusted; repair, do not round}

Every implemented search here follows the same shape, and the shape is the
point.

\begin{lstlisting}
searchCone(problem, degree, budgets):
  columns := boundedSquaresTimesConstraintProducts(problem)
            ++ signedIdealColumns(problem.equalities)
  support := floatingPointLP(columns, target)     -- untrusted
  if no support: return UNKNOWN
  coefficients := exactRationalSolve(support, target)
  if unavailable: return UNKNOWN
  certificate := assemble(coefficients)
  if exactCheck(problem, certificate): return certificate
  return UNKNOWN
\end{lstlisting}

A floating result is not a certificate. The prototypes take the support of the
proposed coefficient vector, re-solve the corresponding coefficient equations
over exact rationals, require the repaired coefficients to be nonnegative, and
expand the certificate to check the original identity. Rationalisation uses a
denominator cap; tiny weights may be dropped during \emph{search
reconstruction}, which never weakens the final check.

This separation handles the most important failure mode: a numerical solver can
treat a small nonzero residual as acceptable, whereas an identity proof cannot.
A rational coefficient merely close to a valid one is still wrong, and a residual
of $10^{-12}$ is not a proof. The test suites include a targeted attack on this
--- a weight of $(10^{20}+1)/10^{20}$ offered where $1$ is required --- which is
why the checkers assert the weight's exact rational type before comparing it.

Numerical infeasibility is likewise not a proof that the goal is false, nor even
a formally established proof that the finite cone contains no solution. Indeed
the implementations are \emph{not complete even for every feasible instance of
their own finite LP cone}, because the support and free-parameter heuristics can
miss an exact feasible solution. And even a mathematically exact LP
infeasibility certificate would refute \emph{that coefficient representation
problem}, not the original theorem. The
repair procedures are deliberately modest: they take coefficients above a fixed
support threshold and set free parameters of the rational system to zero, which
can miss valid certificates with small coefficients or a different feasible
choice of free parameters. The correct result is \unknown{}, not a weakened
acceptance test. A better oracle can explore nearby supports and solve exact
feasibility in the nullspace without changing the checker.

\subsection{Building the dictionary}

For each square-free subset or multiset $I$ of the nonnegative premises with
$|I|\le s$, form the product $g_I$ and enumerate monomials of degree at most
$\lfloor(D-\deg g_I)/2\rfloor$. Take squares of those monomials and of their
pairwise sums and differences. Add signed monomial multiples of each equality
as ideal columns.

Two refinements repay their cost. The first is \emph{target-driven} square
proposal: if the target contains $a u^2+cuv+bv^2$ with $a>0$, propose
$q=u+(c/2a)v$, and when $b>0$ also the opposite orientation; likewise propose
binomial roots from exact rational square-root ratios of positive even-degree
terms. One implementation's initial heuristic recognised $(x-2/3)^2$ but was
defeated by $(x-2/3)^2+1/97$, where a positive constant obscured the coefficient
ratio; adding the cross-coefficient construction recovered the intended square,
and the case is now a regression test.

The second is \emph{variable-support partitioning}: partition the constraints by
variable support and request cross-component products only when the target
connects those components, instead of expanding a global Cartesian product.
Recorded generator budgets run to a few thousand columns; with $M$ monomials an
unguarded binomial dictionary already has up to $M+6\binom{M}{2}$ entries before
deduplication.

\subsubsection*{Four honest failures of a finite dictionary}

The prototypes each kept a control showing a true statement their own search
misses. They are different arguments and all four are worth keeping.

\begin{itemize}[leftmargin=1.4em]
\item \textbf{Motzkin.} $M(x,y)=x^4y^2+x^2y^4+1-3x^2y^2$ is nonnegative by
  arithmetic--geometric mean on its three nonnegative terms, and is famously not
  a sum of squares of polynomials. Bounded square dictionaries return \unknown{}.
\item \textbf{$(2x-3y)^2$, with a counting argument.} A dictionary restricted to
  $m$, $m+n$ and $m-n$ cannot represent a square with unequal coefficients:
  covering the $-12xy$ cross term using $(x-y)^2$ requires weight $6$, already
  exceeding the target's $x^2$ coefficient $4$, and adding other positive
  dictionary atoms cannot repair the deficit.
\item \textbf{$(x+y+1)^2$, with a rank argument.} A rank-one quadratic Gram
  matrix for the coefficient vector $(1,1,1)$ cannot be decomposed as a positive
  sum of rank-one vectors with only two nonzero coordinates unless the summands
  lie in its one-dimensional range, and none of the allowed binomial vectors
  does.
\item \textbf{$(2x+y)^2$ at degree cap 2.} A restricted coefficient-ratio
  dictionary misses it outright --- a goal that the exact quadratic decomposition
  of \S\ref{sec:quadratic} decides immediately.
\end{itemize}

The last of these is the most useful, because it is a \emph{cross-engine}
argument rather than a limitation. One prototype's dictionary obstruction and
another's quadratic completeness theorem together justify the multi-engine
design better than either does alone. A failure of one search is a reason to
route, not a reason to stop.

\begin{remark}
A nonnegative polynomial need not be a sum of polynomial squares at all. Failure
of a bounded search proves only that this search found nothing.
\end{remark}

\subsection{Exact rational quadratic decomposition}
\label{sec:quadratic}
\status{Implemented, tested, and used to generate Lean replay candidates.}

For a rational quadratic polynomial in $n$ variables let
$z=(1,x_1,\ldots,x_n)^{T}$ and $p(\bm x)=z^{T}Qz$, where $Q$ is the unique
symmetric matrix obtained by placing constant and square coefficients on the
diagonal and half of each mixed or linear coefficient in the corresponding
symmetric off-diagonal entries.

\begin{theorem}[Quadratic fragment]
\label{thm:quadratic}
A rational quadratic polynomial is nonnegative on all of $\R^n$ if and only if
its homogenised symmetric matrix $Q$ is positive semidefinite. In that case it
admits a certificate $p=\sum_k d_k\ell_k^2$ with nonnegative rational $d_k$ and
rational affine linear forms $\ell_k$.
\end{theorem}
\begin{proof}
If $Q$ is positive semidefinite the first assertion is immediate. Conversely,
for a vector $(t,\bm y)$ with $t\ne0$,
\[
 (t,\bm y)^{T}Q(t,\bm y)=t^2p(\bm y/t)\ge0,
\]
and the assertion for $t=0$ follows by continuity; thus $Q$ is positive
semidefinite.

For the rational decomposition, choose a positive diagonal pivot $d=Q_{kk}$.
Completing the square gives
\[
 z^{T}Qz=d\left(z_k+\sum_{j\ne k}\frac{Q_{kj}}{d}z_j\right)^2
    +z_R^{T}\!\left(Q_{RR}-\frac{Q_{Rk}Q_{kR}}{d}\right)z_R .
\]
The residual Schur complement is positive semidefinite and every entry remains
rational. Repeat. If all remaining diagonal entries vanish in a positive
semidefinite matrix, all remaining off-diagonal entries also vanish: otherwise a
suitable vector supported on two coordinates gives a negative value. The process
therefore terminates with the claimed rational weighted squares.
\end{proof}

\begin{lstlisting}
quadratic_certificate(p):
  if total_degree(p) > 2: return Unknown
  Q := exact_homogenized_matrix(p)
  active := all coordinates
  terms := []
  while active is not empty:
    if any active diagonal is negative: return Unknown
    choose a positive diagonal pivot d
    if no pivot exists:
      if the active submatrix is zero: break
      else: return Unknown
    emit d * (pivot_variable + rational_linear_tail)^2
    Q := exact Schur complement on the remaining coordinates
  certificate := sum of emitted terms
  return certificate only if expand(certificate) == p
\end{lstlisting}

The algorithm rejects a negative diagonal pivot or a nonzero zero-diagonal
residual, and the returned certificate is checked again by polynomial expansion
through the \emph{general} cone checker. No numerical eigenvalue threshold is
used: a matrix with a tiny nonzero rational off-diagonal entry is treated
exactly, not rounded into positive semidefiniteness.

The matrix arithmetic uses $O(n^3)$ rational operations, which is not an
$O(n^3)$ bit-complexity claim --- numerators and denominators can grow
substantially, and sparse pivots, fraction-free variants and explicit bit-length
budgets are natural production improvements. A fraction-preserving
$LDL^{\mathsf T}$ decomposition is the same algorithm in matrix form.

A failed global quadratic test does not imply the constrained goal is false: $x$
is not globally nonnegative but is nonnegative under the hypothesis $x\ge0$.
Such a result routes to a constraint-aware worker; it does not terminate the
tactic.

Expanding a nonnegative quadratic hides the linear forms whose squares establish
its sign. Existing arithmetic preprocessing may discover the right consequences,
but it need not start from those forms. This worker reconstructs the certificate
systematically rather than hoping a useful square already appears syntactically.

\subsection{Discovered certificates worth recording}

The implemented searches returned, among others:
\begin{align*}
 (a^2+b^2)(c^2+d^2)-(ac+bd)^2&=(ad-bc)^2,\\
 x^2+y^2+z^2-xy-yz-zx
   &=\tfrac12(x-y)^2+\tfrac12(y-z)^2+\tfrac12(z-x)^2,\\
 x^4+y^4-2x^2y^2&=(x^2-y^2)^2,\\
 x^4+y^4+1-x^2y^2-x^2-y^2
   &=\tfrac12(x^2-y^2)^2+\tfrac12(x^2-1)^2+\tfrac12(y^2-1)^2,\\
 x^4+y^4+z^4-x^2y^2-y^2z^2-z^2x^2
   &=\tfrac12(x^2-y^2)^2+\tfrac12(y^2-z^2)^2+\tfrac12(z^2-x^2)^2,\\
 (x^2+y^2+z^2)^2-3(x^2y^2+y^2z^2+z^2x^2)
   &=\tfrac12\bigl((x^2-y^2)^2+(x^2-z^2)^2+(y^2-z^2)^2\bigr),\\
 x-x^3&=x^2(1-x)+x(1-x),\\
 x^2+y^2-2&=(x-y)^2+2(xy-1),\\
 x^3+y^3-xy(x+y)&=(x+y)(x-y)^2\quad(x,y\ge0),
\end{align*}
together with the manually supplied discriminant identity
\begin{equation}\label{eq:disc}
 b^2-4ac=(2ad+b)^2-4a(ad^2+bd+c),
\end{equation}
which is accepted by the checker although it lies outside every implemented
search dictionary --- a deliberate demonstration that the checker's language is
larger than the searcher's.

A domain-dependent example shows why hypothesis products matter. Under
$0\le x\le1$ and $0\le y\le1$,
\[
 2x(1-x)+3y(1-y)+xy+\tfrac15\ge0 .
\]
Supplied in expanded form, with negative $x^2$ and $y^2$ coefficients, the
useful certificate consists of products of the four bound polynomials
$x,1-x,y,1-y$. Global SOS search alone cannot establish this, because the
assertion is domain-dependent. Removing or changing a hypothesis makes the exact
checker reject the certificate, so hypothesis provenance is part of the tested
interface.

\subsection{A concrete path to general SOS search}

A later worker can replace the fixed square list with Gram matrices: choose a
monomial vector $v_D$ and seek $Q_\alpha\succeq0$ satisfying coefficient
identities for
\[
 p= v_D^{T}Q_0v_D+\sum_\alpha v_{D_\alpha}^{T}Q_\alpha v_{D_\alpha}g^\alpha
     +\sum_j u_jh_j .
\]
CVXPY or CSDP can serve as modelling front ends~\cite{cvxpy,csdp}. Their output
remains untrusted: recover rational matrices inside the exact affine
coefficient-constraint space, perform an exact positive-semidefiniteness check,
and convert each matrix to rational weighted squares. Entrywise rounding
followed by approximate eigenvalue testing is not acceptable. Singular Gram
matrices are especially delicate, since perturbation can destroy either
positivity or the coefficient identity; exact nullspace handling, support
changes, and facial reduction are worthwhile \emph{search} techniques, but
acceptance must still reduce to \eqref{eq:cone}~\cite{rational-sos}. No
implementation here contains an SDP solver and no SDP performance claim is made.

\begin{remark}[Reuse the existing Lean backend]
\label{rem:sos-reuse}
This point deserves emphasis because it materially narrows what \forge{} should
claim. A Lean \code{sos} project already implements an end-to-end
nonlinear-real-arithmetic tactic after Harrison's decision procedure: reification,
search, rational reconstruction, and certificate-based proof construction, with
a witness-oriented replay interface and a configuration that already includes
degree deepening, constraint-product bounds, and optional refutation
powers~\cite{lean-sos,sos-engine,sos-core}. Consequently ``add an SOS tactic to
Lean'' is \emph{not} a novel component of this design. Similarly, a Lean
\code{lp} project supplies verified linear programming and a supported
quantified rational-affine fragment~\cite{lean-lp,lp-blog}, covering much of the
affine-witness machinery.

The production arrangement should therefore be staged: run the cheap in-house
certificate searches first, and escalate to the existing libraries rather than
building a second unaudited wrapper. \forge{}'s contribution in this area is
\emph{when} to request a certificate, \emph{how} to bind it to the source
problem, and \emph{what} to do with the proved consequence afterwards --- not
the certificate search itself. The version caveat of \S\ref{sec:baseline}
applies: compatibility between those projects and the pinned baseline must be
established, not assumed.
\end{remark}

\subsection{Bernstein certificates on rational boxes}
\status{Implemented and tested in Python by seven independent efforts.}

\subsubsection*{Coefficients and the bound}

Let $p$ be a polynomial on the rational box
$X=\prod_{i=1}^{n}[l_i,u_i]$ with $l_i<u_i$. Substituting
$x_i=l_i+(u_i-l_i)t_i$ gives $q(\bm t)=\sum_\alpha a_\alpha\bm t^\alpha$ on the
unit box. With $d_i$ a coordinatewise degree bound, the tensor-product Bernstein
expansion is
\[
 q(\bm t)=\sum_{0\le\beta\le\bm d} b_\beta
   \prod_i\binom{d_i}{\beta_i}t_i^{\beta_i}(1-t_i)^{d_i-\beta_i},
\]
with exact rational coefficients
\begin{equation}
 b_\beta=\sum_{\alpha\le\beta}a_\alpha
     \prod_i\frac{\binom{\beta_i}{\alpha_i}}{\binom{d_i}{\alpha_i}} .
 \label{eq:bernstein}
\end{equation}

\begin{lemma}
For $0\le a\le d$,
$\displaystyle t^a=\sum_{b=a}^{d}\frac{\binom ba}{\binom da}\binom db\,
t^b(1-t)^{d-b}$.
\end{lemma}
\begin{proof}
Use $\binom db\binom ba=\binom da\binom{d-a}{b-a}$, factor out $t^a$, and apply
the binomial theorem to the remaining sum, which collapses to
$(t+(1-t))^{d-a}=1$.
\end{proof}

Applying this in each coordinate proves \eqref{eq:bernstein}. Each basis term is
nonnegative on the unit box and the basis terms sum to one, so
\[
 \min_\beta b_\beta\le p(\bm x)\le\max_\beta b_\beta\qquad(\bm x\in X).
\]
In particular nonnegative coefficients certify nonnegativity throughout the box,
and strictly positive coefficients certify strict positivity.

\begin{remark}
Negative Bernstein coefficients are \emph{not} counterexamples. They mean only
that this sufficient condition has not proved the desired sign. A returned
negative point is a different kind of result and must be evaluated exactly.
\end{remark}

\subsubsection*{The subdivision tree is part of the proof}

If a box does not certify immediately, split a side at an interior rational
point. A certificate is a binary tree: an internal node contains an axis, an
interior rational split point, and \emph{both} children; a leaf contains its
coefficients or a claimed nonnegative rational lower bound.

The checker independently reconstructs every child box from the root, recomputes
each leaf's coefficients, and confirms the claimed bound. Requiring both
children proves coverage --- there is no way to omit a difficult corner by
supplying an incomplete list of favourable boxes. A missing child, an invalid
split axis, or a split point at an interval endpoint is rejected, and every
basis index must appear exactly once in canonical order. Overlap on a shared
face is harmless; missing an entire child is not. Node and depth limits are
operational rejection limits, not mathematical assumptions about positivity.

The strongest implementations make the checker algorithmically \emph{opposite}
to the producer: rather than recomputing \eqref{eq:bernstein}, they expand the
claimed Bernstein representation --- in the original $x$ coordinates, using
\[
 \binom{d}{k}\frac{(x-a)^k(b-x)^{d-k}}{(b-a)^d},
\]
--- and compare the result against an independently computed affine transform of
the target. Producer and checker then share no conversion formula, though they
do still share the sparse rational arithmetic layer, so their independence is
not complete.

Search may evaluate the polynomial at a midpoint to abandon a hopeless box
early, but it must never return that evaluation as a refutation. Some
implementations do return a checked refutation: an exact point inside the box at
which the polynomial is negative, validated independently for membership and
sign. That is a genuine three-valued outcome ---
\textsc{proved}/\textsc{refuted}/\unknown{} --- and it is the only sound negative
route any of the prototypes implements.

Split-axis policy matters for the completeness argument below. Cycling axes as
$\mathrm{depth}\bmod n$, or choosing the largest \emph{relative} width, both
ensure every coordinate shrinks along a branch; choosing the largest absolute
width introduces a coordinate-unit bias.

\subsubsection*{A qualified completeness result, and a sharp obstruction}

\begin{proposition}[Strictly positive polynomials]
\label{prop:bernstein-complete}
For a polynomial strictly positive on a compact box, a fair subdivision search
that shrinks every coordinate along each branch eventually makes all Bernstein
coefficients positive.
\end{proposition}
\begin{proof}[Justification]
The transformed higher-degree coefficients tend uniformly to zero as box widths
shrink, while the constant terms stay above a positive minimum: each local
Bernstein coefficient is $p(a)+O(\max_i(b_i-a_i))$ with a bound uniform over the
original box. Compactness supplies the positive minimum.
\end{proof}

This does not give a useful bound for every input and does not justify ignoring
the configured caps.

Zeros behave completely differently, and the difference is a theorem rather than
a matter of search depth.

\begin{theorem}[Interior-zero obstruction]
\label{thm:zero}
Let $p$ be a nonzero polynomial with a zero in the interior of a box $B$. Then
$p$ has no all-nonnegative Bernstein expansion on $B$.
\end{theorem}
\begin{proof}
Every Bernstein basis function is strictly positive at an interior point. If all
coefficients were nonnegative, a zero value at an interior point would force
every coefficient to vanish, hence $p\equiv0$ on $B$ and so $p\equiv0$.
\end{proof}

Consequently a finite dyadic subdivision cannot certify $(x-1/3)^2$ on $[-1,1]$
by nonnegative leaf coefficients alone: some leaf must contain $1/3$ in its
interior. The obstruction is quantitative. On $[0,1]$, the dyadic cell
containing $1/3$ always has the root in its interior, and if its endpoints are
$\ell<u$ then the middle quadratic Bernstein coefficient is exactly
\[
 (\ell-\tfrac13)(u-\tfrac13)<0 .
\]
No finite midpoint-only subdivision can certify that cell by nonnegative
degree-two coefficients, whereas the exact quadratic engine of
\S
ef{sec:quadratic} proves the same inequality immediately. This is a
concrete reason to cooperate across representations rather than escalate one
representation's resource cap indefinitely. The same obstruction applies with an \emph{irrational} zero, and there
it rules out any rational-box tree whatever: $(x^2-2)^2$ on $[1,2]$ has its zero
at $\sqrt2$, which no rational split point can ever fall on. The prototypes
record both as expected \unknown{} outcomes, at depth 12 and at their depth
limits respectively.

A second, milder limitation is worth recording: $(x-y)^2$ on $[0,1]^2$ returns
\unknown{} at depth four even though a single square certificate proves it
immediately. Again this routes rather than stops.

\subsubsection*{Worked instances}

For $p(x)=(x-1/2)^2+1/100$ on $[0,1]$, the degree-two coefficients at the root
are $13/50,\,-6/25,\,13/50$, so the root check cannot certify positivity.
Splitting at $1/2$ gives
\[
 (13/50,\,1/100,\,1/100)\quad\text{and}\quad(1/100,\,1/100,\,13/50),
\]
so two leaves suffice for strict positivity. For $p(x)=x^2-x+1/3$ on $[0,1]$ the
root coefficients are $(1/3,-1/6,1/3)$, and splitting at $1/2$ gives
$(1/3,1/12,1/12)$ and $(1/12,1/12,1/3)$. For $p(x)=x^2-x+3/10$ the split yields
two leaves each with exact lower bound $1/20>0$ and a three-node certificate.
These are exact calculations, not numerical plots.

A multivariate instance uses the bounded Motzkin polynomial
$M(x,y)=x^4y^2+x^2y^4-3x^2y^2+1$. The search certifies $M+1/16$ on three
domains: on $[-1,1]^2$ a single leaf suffices, with minimum coefficient exactly
$1/16$; on $[-2,2]^2$ it constructs 32 leaves at maximum depth six; on
$[-3,3]^2$, 56 leaves at maximum depth nine.

\begin{remark}[An instructive correction]
In one development run the first test wrongly expected the smallest domain to
require subdivision. The exact coefficient calculation showed the test
assumption was wrong. The corruption test was moved to the larger domain and
independent symbolic reconstruction tests were added. No positivity condition
was weakened to make the test pass. That is the correct response to a failing
test in this setting, and it is recorded here because the opposite response is
the standard way such a project becomes unsound.
\end{remark}

\subsubsection*{Cost}

The tensor has $\prod_i(d_i+1)$ coefficients and a binary tree can have
exponentially many leaves in its depth bound; exact coefficient bit sizes also
grow under repeated substitution. The method suits moderate dimension, low
coordinatewise degree, and meaningful compact bounds. Implementations cap tensor
size before expansion precisely because an adversarial certificate could
otherwise request an enormous one. This is attractive research code, not a
complete defence against hostile memory-exhaustion input, and not a universal
substitute for nonlinear arithmetic.

\subsection{Univariate certificates that preserve zeros}
\label{sec:sturm}
\status{Implemented and tested in Python by one effort. No generic Lean Sturm
replayer exists.}

Theorem~\ref{thm:zero} is not merely a caveat; it is a capability gap with a
constructive answer. This worker is the answer, and it is the only mechanism in
the whole collection that can certify a polynomial with an interior zero.

\subsubsection*{Square-factor reduction}

For a nonzero rational univariate polynomial, extract even multiplicities:
\[
  p(x)=q(x)^2\,r(x),
\]
where $r$ is square-free and its rational scalar coefficient retains the
original sign. Full factorisation is not necessary in principle --- square-free
decomposition suffices. The prototype uses a rational factorisation interface as
an \emph{oracle} and then verifies the resulting identity with its own exact
representation~\cite{sympy-polys}.

A certificate on a rational interval $[a,b]$, $a<b$, consists of $q$, $r$, and a
Sturm chain for $r$. The checker requires:
\begin{enumerate}[leftmargin=1.4em]
\item the exact identity $p=q^2r$, and a valid terminating signed-remainder
  chain for $r$;
\item no roots of $r$ in the \emph{open} interval $(a,b)$;
\item $r((a+b)/2)>0$ and $p(a),p(b)\ge0$.
\end{enumerate}
The zero polynomial is handled by $q=0$, $r=1$.

\subsubsection*{The checked Sturm chain}

The chain starts with $r_0=r$, $r_1=r'$, and continues with
$r_{i+1}=-\operatorname{rem}(r_{i-1},r_i)$ until the next remainder is zero; a
constant $r$ has a one-element chain. The checker recomputes the derivatives and
rational remainders, requires exact equality with the supplied chain, and
requires the last element to be a nonzero constant --- which establishes that $r$
is square-free. Crucially it \emph{recomputes the entire chain} rather than
trusting a stored root count.

Let $V(t)$ denote the number of sign changes in the chain. Endpoint roots are
the easy thing to get wrong here, so the implementation computes \emph{one-sided}
signs: at $a+$, the sign of a chain polynomial is the sign of its first nonzero
derivative at $a$; at $b-$, that sign is multiplied by $(-1)^k$ when the first
nonzero derivative has order $k$. The open-interval root count is then
\[
 N_{(a,b)}(r)=V(a+)-V(b-).
\]
Independent tests compare root counts against a symbolic oracle and explicitly
exercise rational roots at both endpoints.

\begin{theorem}[Univariate interval certificate]
A certificate satisfying the three conditions above proves $p(x)\ge0$ for every
$x\in[a,b]$.
\end{theorem}
\begin{proof}
The Sturm count establishes that $r$ has no root in $(a,b)$. By continuity its
sign is constant there, and the positive midpoint value fixes that sign as
positive. Hence $q^2r\ge0$ on the interior. The exact identity and the two
endpoint checks complete the proof on the closed interval.
\end{proof}

\begin{proposition}[Fragment completeness]
Ignoring degree and resource caps, exact square-free decomposition together with
this check is complete for nonnegativity of rational univariate polynomials on a
rational closed interval.
\end{proposition}
\begin{proof}
If a nonzero $p=q^2r$ is nonnegative, a simple interior root of the square-free
$r$ would change the sign of $p$, which is impossible. Thus $r$ has no interior
root and is positive there, and the certificate theorem applies.
\end{proof}

This argument uses the mathematical Sturm theorem. Its use is justified in the
algorithm specification and exercised by tests, but it is \emph{not} formalised
in Lean anywhere in this collection. A production Lean implementation must prove
or integrate that theorem and check the instantiated certificate before this
worker can be trusted; that is a substantial verification task and must not be
hidden behind an unchecked foreign call.

\subsubsection*{Why it belongs alongside Bernstein search}

$(x-1/3)^2$ is handled without aligning a dyadic split with its zero, and so is
$(x^2-2)^2$, whose real zeros are irrational. A constructed test family
multiplies rational and irrational even-root factors by a positive quadratic and
the endpoint factors $(x+2)(2-x)$ on $[-2,2]$. All 30 instances are certified by
the square-factor/Sturm worker; \emph{none} is certified by the depth-12
Bernstein-only ablation.

That is a genuine component-level distinction. It does not imply these examples
are difficult for an existing Lean tactic once a factorisation is supplied. The
contribution is an explicit search-and-check route for \emph{finding} and
\emph{validating} the relevant algebraic structure.

In several variables, candidate factorisation, algebraic equality reasoning, and
constraint-relative certificates can play an analogous role. A generic
multivariate zero-set treatment is not part of any prototype.

\subsection{Cheap facts for other engines: certified product envelopes}
\status{Specified; not separately implemented.}

Bound propagation can generate inexpensive facts before any larger nonlinear
search runs. If $l_x\le x\le u_x$, $l_y\le y\le u_y$, and $z=xy$, then the four
products of nonnegative bound differences give
\begin{align*}
 z&\ge l_xy+l_yx-l_xl_y, & z&\ge u_xy+u_yx-u_xu_y,\\
 z&\le u_xy+l_yx-u_xl_y, & z&\le l_xy+u_yx-l_xu_y.
\end{align*}
The first, for instance, is the expansion of $(x-l_x)(y-l_y)\ge0$. These facts
are \emph{linear} in $x,y,z$ and can be consumed immediately by the existing
arithmetic engine; their derivations, bounds, and the equation $z=xy$ must
accompany them.

This worker should be prioritised because its certificates are small and its
output is useful across theories, and it has a natural trigger: run only when
the bounds of a product operand improve. Conversely, an extracted alternative
form such as $(x-y)^2+2(xy-1)$ can become a positivity candidate for
$x^2+y^2-2$, and a cone certificate can request a particular difference term
$x-y$ as a guide for other matching tasks. What is shared between views is a
\emph{proved relation}, never a mutable polynomial whose meaning changes during
search.

\subsection{A complete cross-theory proof shape}

Suppose $p=x^2-2xy+y^2$ with hypothesis $p\le0$, and the target is $f(p)=f(0)$
for an otherwise uninterpreted $f:\R\to\R$. The nonlinear worker discovers
$p=(x-y)^2$, proves $0\le p$, and exports that inequality; antisymmetry gives
$p=0$; congruence closure gives the target. No worker needs a theory of $f$, and
no single worker solves the whole problem.

This is the central integration pattern of the whole design. It is not an
assertion that stock \grind{} fails on this small example.

\subsection{Domains, casts, and division}

The controller must establish that the source assumptions imply the box bounds
before applying a Bernstein theorem, and polynomial reification must preserve
variable identity, source operations, and casts. A real-field certificate cannot
be applied to an arbitrary ring or a field of finite characteristic.

For natural or integer inequalities, prove the embedding into an ordered field
explicitly and translate the relevant operations. Integer division, truncated
subtraction, and modular arithmetic are not field operations. Translating
\code{(a - b : Nat)} to a field difference requires a guard such as $b\le a$ and
an appropriate transport theorem; otherwise the operation stays opaque or the
relevant cases are split.

Unknown real-valued applications can be polynomial atoms without being
uninterpreted everywhere. One engine may know a theorem about $f(x)$ while the
polynomial engine treats its value as a single coordinate; once that theorem
yields a bound or equality, the bound is proved in the original context and can
then be translated into the polynomial representation. No global erasure of
dependent types is necessary.

Finally, if congruence closure identifies two source terms, a worker may share
one polynomial variable only after preserving the equality proof that justifies
the identification. It must not merge different variables because their printed
names agree. Treating equal terms as distinct opaque variables is a loss of
proving power rather than an unsound strengthening --- which is precisely why the
shared-state integration is worth the engineering.
